% ---------------------------------------------------------------------------
%  中文简历：一页 A4
%  左栏时间、右栏经历；图标用内置的 fontawesome5，中文字体用系统自带的。
%  想换中文字体：\setCJKmainfont{字体名}（本机没有的字体，首次编译时日志里会提示）。
%  人物与经历均为示例。写简历的原则：每一条都回答「做了什么、结果如何」。
% ---------------------------------------------------------------------------
\documentclass[UTF8,zihao=5]{ctexart}

\usepackage[a4paper,top=1.5cm,bottom=1.3cm,left=1.7cm,right=1.7cm]{geometry}
\setmainfont{texgyrepagella}[
  Extension=.otf, UprightFont=*-regular, ItalicFont=*-italic,
  BoldFont=*-bold, BoldItalicFont=*-bolditalic]
\usepackage{fontawesome5}
\usepackage{xcolor}
\usepackage{enumitem}
\usepackage{tabularx}
\usepackage[hidelinks]{hyperref}

% 墨色做结构，绿色是唯一的强调
\definecolor{ink}{HTML}{1F3A33}
\definecolor{accent}{HTML}{2F7D63}
\definecolor{mist}{HTML}{75837E}

\pagestyle{empty}
\linespread{1.38}
\setlength{\parindent}{0pt}
\setlength{\parskip}{0pt}
\setlist[itemize]{nosep, leftmargin=1em, label={\color{accent}\rule[.3ex]{3pt}{3pt}}}

\newlength{\datecol}
\setlength{\datecol}{6.2em}

% 栏目：\block{教育背景}，左栏写标题，右边一条细线
\newcommand{\block}[1]{%
  \par\vspace{14pt}
  \makebox[\datecol][l]{\color{accent}\bfseries #1}%
  {\color{mist!40}\leaders\hrule height .4pt\hfill}\par\vspace{6pt}}

% 经历：\entry{时间}{单位}{角色}{一句话}，一句话可留空；下面接 body 里的 itemize
\newcommand{\entry}[4]{%
  \par\noindent
  \parbox[t]{\datecol}{\small\color{mist}#1}%
  \parbox[t]{\dimexpr\linewidth-\datecol}{%
    {\bfseries\color{ink}#2}\hfill{\small\color{mist}#3}\par
    \ifx\relax#4\relax\else{\small\color{ink!80}#4}\par\fi}\par\vspace{2pt}}

% 正文缩进到右栏
\newenvironment{body}
  {\begin{list}{}{\setlength{\leftmargin}{\datecol}\setlength{\topsep}{0pt}\setlength{\parsep}{0pt}}\item\relax}
  {\end{list}\vspace{6pt}}

\begin{document}

% ---------------------------------------------------------------- 抬头
{\fontsize{28pt}{32pt}\selectfont\bfseries\color{ink}林知远}%
\hspace{1em}{\large\color{mist}设计工程师 · 阅读与排版产品}\par
\vspace{8pt}
{\small\color{mist}
  \faPhone*\enspace 138-0000-0000\qquad
  \faEnvelope[regular]\enspace\href{mailto:linzhiyuan@example.com}{linzhiyuan@example.com}\qquad
  \faGithub\enspace\href{https://github.com/example}{github.com/example}\qquad
  \faMapMarker*\enspace 杭州\par}
\vspace{10pt}
\begin{body}
  六年做一件事：让屏幕上的文字更好读。既写界面代码，也画版式；
  相信好的工具应当安静，把注意力还给内容本身。
\end{body}

% ---------------------------------------------------------------- 经历
\block{工作经历}

\entry{2022.07\\至今}{某某阅读}{高级设计工程师 · 排版组负责人}%
  {负责阅读器的排版引擎与版式系统，服务 3\,000 万月活读者}
\begin{body}
\begin{itemize}
  \item 重写中文断行与标点挤压模块，长段落排版耗时从 48\,ms 降到 6\,ms，翻页卡顿投诉下降 72\%；
  \item 设计「舒适版式」：依据行长与行距实验确定默认参数，读者平均单次阅读时长提升 11\%；
  \item 建立设计与工程共用的版式令牌（字号、行距、留白），新页面开发周期从两周缩短到四天。
\end{itemize}
\end{body}

\entry{2020.03\\2022.06}{某某科技}{前端工程师}%
  {负责知识库产品的编辑器与文档渲染}
\begin{body}
\begin{itemize}
  \item 实现所见即所得的公式与表格编辑，覆盖 95\% 的 \LaTeX{} 常用语法；
  \item 主导编辑器性能治理，万字文档的输入延迟稳定在 16\,ms 以内。
\end{itemize}
\end{body}

% ---------------------------------------------------------------- 项目
\block{个人项目}

\entry{2021 -- 今}{「行间」中文排版工具集}{作者 · 开源}{}
\begin{body}
\begin{itemize}
  \item 一组处理中西文混排间距、标点禁则与避头尾的小工具，GitHub 2.4k Star，被 30 余个项目采用；
  \item 撰写系列文章《中文排版的十个细节》，累计阅读 50 万次。
\end{itemize}
\end{body}

% ---------------------------------------------------------------- 教育
\block{教育背景}

\entry{2016 -- 2020}{某某大学}{数字媒体技术 · 本科}%
  {GPA 3.7/4.0；毕业设计《屏幕中文正文的最适行长研究》获校级优秀论文}

% ---------------------------------------------------------------- 技能
\block{专业技能}

\begin{body}
\renewcommand{\arraystretch}{1.25}
\begin{tabularx}{\linewidth}{@{}>{\color{mist}}l X@{}}
  工程 & TypeScript、React、Rust、WebAssembly、Canvas 与文本渲染 \\
  设计 & 版式与字体设计、设计系统、Figma、交互原型 \\
  排版 & \LaTeX{}、OpenType 特性、Unicode 断行规则、W3C《中文排版需求》 \\
  语言 & 英语（工作语言，CET-6 620）、日语（N2） \\
\end{tabularx}
\end{body}

\end{document}
